\section{Upper volume deficit and balanced Weyl moments}
\label{sec:volume}

Write $V_M=\operatorname{Vol}(M,g)$ and $s_0=\mathbf E(v^2+b^2)$.
Our goal in this section is to prove that the two second moments agree
whenever neither Weyl block vanishes. The characteristic identities
turn an upper polar-volume deficit into an integer restriction; a
weighted scalar variance estimate then excludes the sole remaining
unequal-moment case.

Radial comparison also plays a central role in Cheng's classification
under the additional assumption $\chi(M)\le3$
\cite[Proposition~3.1]{Cheng2026}. His Riccati argument does not require
curvature matrices at different times to commute. We use a discrete
Jacobi determinant to retain their order and obtain the following
moment-dependent deficit.

\begin{proposition}[Upper volume deficit]\label{prop:upper-volume-deficit}
One has
\[
 \frac{V_M}{8\pi^2/3}\le 1-\frac{3}{40}s_0.
 \tag{7}\label{eq:7}
\]
\end{proposition}

Here is how the estimate will be used. If both blocks are nonzero, their
gaps and $v+b<2$ give $2\le s_0<4$. Equations~\eqref{eq:6} and
\eqref{eq:7} then imply $\chi<7$ and $|\tau|/\chi<4/15$.
The parity and nonnegativity constraints from $n_\pm$ leave only
$(\chi,|\tau|)=(5,1)$ when $\tau\ne0$. The last part of this section
rules out that case and proves Proposition~\ref{prop:balanced-moments}.
We first establish the upper deficit and the common radial machinery;
the conditional lower comparison needed after one-block vanishing is
proved in Section~\ref{sec:classification}.

Along a unit-speed geodesic, choose a parallel orthonormal normal frame
and let $\mathcal R(t)$ be the normal Jacobi curvature matrix.  Set
\[
 A(t)=\mathcal R(t)-\mathrm{Id},\qquad
 \mathcal J''+\mathcal R\mathcal J=0,\qquad
 \mathcal J(0)=0,\quad \mathcal J'(0)=\mathrm{Id},\qquad
 D(t)=\det\mathcal J(t).
\]
Positive sectional curvature and $\operatorname{Ric}=3g$ give
$0<\mathcal R<3\mathrm{Id}$ and $\operatorname{tr}A=0$.
Let $\mathcal E$ denote integration over the unit tangent bundle against
normalized Riemannian measure on the base and normalized spherical
measure in the fibers.

\begin{lemma}[Geodesic averages]\label{lem:geodesic-averages}
At every time along the geodesic flow,
\[
 \mathcal E|A|^2=\frac32s_0,\qquad
 \mathcal E|A'|^2\le\frac32\,\mathcal E|A|^2.
 \tag{8}\label{eq:8}
\]
\end{lemma}
\begin{proof}
The indicated probability measure is invariant under geodesic flow.
Indeed, in cotangent coordinates the Hamiltonian flow of kinetic energy
preserves canonical volume: its divergence vanishes by equality of the
mixed second derivatives of the Hamiltonian.  Applying this observation
to energy shells gives invariance of the measure on the unit-energy
level.  In orthonormal fiber coordinates this is precisely Riemannian
base volume times spherical measure.

The plus and minus contributions to $A$ are orthogonal conjugates of
$T/2$ and $U/2$, using the normal bases $I_s\dot\gamma$ and
$J_s\dot\gamma$.  The same description applies to their derivatives in
parallel frames.  The mixed contractions have zero spherical average:
they are invariant bilinear pairings between the inequivalent types
$(2,0)$ and $(0,2)$, or between $(5/2,1/2)$ and $(1/2,5/2)$.
Consequently,
\[
 \mathcal E|A|^2=\frac14\mathbf E(|T|^2+|U|^2),\qquad
 \mathcal E|A'|^2
 =\frac1{16}\mathbf E(|\nabla T|^2+|\nabla U|^2).
\]
Here the second factor $1/4$ in the derivative identity is the spherical
average of a squared directional derivative in four dimensions.
Integrating the equation for $\Delta T$ yields
\[
 \mathbf E e=6\mathbf E(\sigma v^3-v^2)\le6\mathbf E v^2,
\]
since $\sigma\le1$ and $v<2$.  The exchanged identity gives the same
bound for $\bar e$.  The stated formulas follow.
\end{proof}

\subsection{Minimizing geodesics and upper polar integration}

The global input is the classical index-form and closed-geodesic argument
of Myers, Synge, and Klingenberg
\cite{Myers1941,Synge1936,Klingenberg1959}; see also
\cite[Theorem~1.5 and Remark~1.2(ii)]{AbreschMeyer1997}.
Its full proof is in Appendix~\ref{app:topology}, under all the standing
oriented, closed, normalized positivity hypotheses.

\begin{lemma}[Diameter, injectivity, and topology]\label{lem:volume-topology}
The diameter of $M$ is at most $\pi$, its injectivity radius is at least
$\pi/\sqrt3$, and $M$ is simply connected. Every minimizing geodesic
segment has no point conjugate to its initial point in its interior.
\end{lemma}

The last assertion allows a conjugate endpoint. It is needed in the
upper polar argument even for segment lengths between $\pi/\sqrt3$ and
$\pi$.

Put $D_c(t)=D(t)$ before the first conjugate time and $D_c(t)=0$ at and
after it.

\begin{lemma}[Upper polar bound]\label{lem:upper-polar}
One has
\[
 V_M\le2\pi^2\int_0^\pi\mathcal E D_c(t)\,dt.
\]
\end{lemma}
\begin{proof}
Fix a center $p$ and take all vectors of length
less than $\pi$ whose radial geodesics have no conjugate point at or
before their endpoints.  This is an open set: the Jacobi determinant
divided by $t^3$, continued as $1$ at zero, has a positive minimum on
each such radial interval.  Every point is the exponential image of a
minimizing vector of length at most $\pi$, by
Lemma~\ref{lem:volume-topology}.  A minimizing vector is in this open
set unless it is critical for $\exp_p$, has length $\pi$, or is zero.

The exponential map is smooth with bounded derivative on the closed
radius-$\pi$ ball.  The area formula with multiplicity states that the
integral of its absolute Jacobian over a measurable set equals the
integral of the number of preimages over its image.  In particular,
critical vectors and the boundary sphere have images of zero
four-dimensional volume: their respective Jacobian integral and domain
volume are zero.  On the open preconjugate set the determinant is
positive.  The area formula therefore bounds $V_M$ by the integral of
$D$ over this set, with multiplicities only increasing that integral.
Nonconjugate cut vectors are included in the set; no identification of
cut points with conjugate points is used.  Polar coordinates give the
claimed bound at each center, since the area of the unit
three-sphere is $2\pi^2$.  Average over centers and use Tonelli's theorem.

\end{proof}

\subsection{Discrete determinant comparisons}

The upper polar bound reduces the volume problem to estimates for a
matrix Jacobi determinant. We approximate the Jacobi equation by a finite
Dirichlet matrix. Its Schur complements and sine-mode compressions
retain the order of matrix multiplication, without requiring simultaneous diagonalization along the geodesic.
The same discrete determinant will be used again for the conditional
lower comparison in Section~\ref{sec:classification}.

\begin{lemma}[The discrete Jacobi determinant]\label{lem:discrete-jacobi}
Fix a smooth symmetric normal curvature matrix on $[0,t]$, put $a=t/N$,
and let $\mathcal L_N$ be the block tridiagonal matrix on nodes
$1,\ldots,N-1$ with diagonal blocks $2\mathrm{Id}-a^2\mathcal R(ia)$
and off-diagonal blocks $-\mathrm{Id}$.  Then
\[
 a^3\det\mathcal L_N\longrightarrow D(t).
\]
If $t$ is before the first conjugate time, $\mathcal L_N$ is positive
definite for all sufficiently large $N$.  The Jacobi convergence is
uniform over families with uniformly bounded coefficients and first
derivatives.
\end{lemma}
\begin{proof}
Define matrices without using inverses by
\[
 Y_0=0,\quad Y_1=\mathrm{Id},\qquad
 Y_{i+1}=(2\mathrm{Id}-a^2\mathcal R(ia))Y_i-Y_{i-1}.
\]
The position $aY_i$ and backward velocity $Y_i-Y_{i-1}$ evolve by
first changing velocity by $-a\mathcal R(ia)(aY_i)$ and then changing
position by $a$ times the new velocity.  On the exact Jacobi solution
one step has error $O(a^2)$, while the update amplifies errors by at most
$1+Ca$.  Iteration proves convergence of position and velocity, uniformly
on the fixed interval.  The constants are uniform under the stated
coefficient bounds.

Suppose first that $t$ is preconjugate.  Away from zero the exact Jacobi
matrix is invertible with uniformly bounded inverse on compact time
intervals, so the converging $aY_i$ are invertible there.  Near zero use
the identity, still without inverses,
\[
 Y_i=i\mathrm{Id}-a^2\sum_{r<i}(i-r)\mathcal R(ra)Y_r.
\]
For sufficiently small fixed $ia$ it gives $\max_{r\le i}|Y_r|\le Ci$
by absorbing the error, and then
$Y_i/i=\mathrm{Id}+O((ia)^2)$.  Thus all the early matrices are invertible
as well.  All inverses needed for elimination have now been obtained
independently of that elimination.

The Schur blocks are
\[
 S_i=Y_{i+1}Y_i^{-1}
 =2\mathrm{Id}-a^2\mathcal R(ia)-S_{i-1}^{-1}.
\]
They are symmetric, starting with
$S_1=2\mathrm{Id}-a^2\mathcal R(a)$.  Near zero they differ from
$((i+1)/i)\mathrm{Id}$ by $O((ia)^2+a^2)$; away from zero, convergence of
position and velocity gives $S_i=\mathrm{Id}+O(a)$.  Choose the small
initial interval first, and then $N$ large.  All Schur blocks are
positive definite.  Block elimination proves positivity of
$\mathcal L_N$ and gives
\[
 \det\mathcal L_N=\prod_{i=1}^{N-1}\det S_i=\det Y_N.
\]
The desired limit follows.  The determinant identity holds for arbitrary
coefficients by continuity, since invertible Schur blocks form a dense
set; convergence therefore gives the limit without the preconjugacy
assumption as well.
\end{proof}

\begin{lemma}[Upper radial comparison]\label{lem:upper-radial}
For $0<t<\pi$, set
\[
 \beta=\frac{t^2}{\pi^2-t^2},\qquad
 B=\frac2t\int_0^t A(s)\sin^2(\pi s/t)\,ds.
\]
Before conjugacy,
\[
 D(t)\le\sin^3t\,\det(\mathrm{Id}-\beta B),
 \qquad \mathrm{Id}-\beta B\succeq0.
\]
At every time $0<t<\pi$,
\[
 1-\frac{D_c(t)}{\sin^3t}
 \ge\frac16\min\{1,\beta^2\max(1,3-\beta)\}\,|B|^2.
\]
\end{lemma}
\begin{proof}
Compare $\mathcal L_N$ with the matrix for $\mathcal R=\mathrm{Id}$,
whose scalar sine-mode eigenvalues are
$2-2\cos(k\pi/N)-a^2$.  They are positive for large $N$ because $t<\pi$.
Normalize $\mathcal L_N$ by the inverse square root of this comparison
matrix.  The resulting positive matrix has trace $3$ on every diagonal
sine-mode block, since $\operatorname{tr}A(ia)=0$.

Split off its first three-dimensional mode and write the normalized
matrix as
$\left(\begin{smallmatrix}P&C\\C^*&Q\end{smallmatrix}\right)$.
Its Schur complement is positive and bounded above by $Q$.  Thus
\[
 \det\begin{pmatrix}P&C\\C^*&Q\end{pmatrix}
 =\det P\,\det(Q-C^*P^{-1}C)
 \le\det P\,\det Q\le\det P.
\]
The last inequality is arithmetic--geometric mean, since
$\operatorname{tr}Q=\dim Q$.  The first-mode compression tends to
$\mathrm{Id}-\beta B$, because its normalized sine vector has entries
$\sqrt{2/N}\sin(\pi i/N)$ and
\[
 \frac{a^2}{2-2\cos(\pi/N)-a^2}\longrightarrow\beta.
\]
Lemma~\ref{lem:discrete-jacobi} gives the determinant ratio limit
$D(t)/\sin^3t$.  Positive compressions have a positive semidefinite
limit, and no inverse of that limiting compression is used.
This proves the first comparison.

The eigenvalues of $B$ lie in $[-1,2]$ and sum to zero, so $|B|^2\le6$.
Before conjugacy the determinant loss is
\[
 1-\det(\mathrm{Id}-\beta B)
 =\frac{\beta^2}{2}|B|^2+\beta^3\det B.
\]
Only $\det B<0$ is unfavorable.  Write the eigenvalues then as
$-r,a_1,a_2$, where $a_1,a_2\ge0$ and $a_1+a_2=r$.
One has
\[
 |\det B|\le\frac r6|B|^2,\qquad r\le1,\qquad\beta r\le2.
\]
The first inequality follows from $(a_1-a_2)^2\ge0$; the last follows
from $\beta a_i\le1$.  These inequalities give the displayed loss bound.
If $\det B\ge0$ the same bound is immediate.  After conjugacy its left
side is $1$, so $|B|^2\le6$ gives it directly.
\end{proof}

\subsection{An exact constant for the upper deficit}

We use one rational sine minorant for both the upper calculation here
and the conditional lower calculation in Section~\ref{sec:classification}.

\begin{lemma}[A rational sine minorant]\label{lem:sine-minorant}
Put
\[
 a_0=\frac{157}{50},\qquad b_0=\frac{22}{7},\qquad
 S(z)=a_0z-\frac{(b_0z)^3}{6}
          +\frac{(a_0z)^5}{120}-\frac{(b_0z)^7}{5040}.
\]
The constants satisfy $a_0<\pi<b_0$, and
$0\le S(z)\le\sin(\pi z)$ for $0\le z\le1/2$. Its coefficients are
\[
 S(z)=\frac{157}{50}z-\frac{5324}{1029}z^3
       +\frac{95388992557}{37500000000}z^5
       -\frac{155897368}{259416045}z^7.
\]
\end{lemma}
\begin{proof}
The elementary bounds $a_0<\pi<b_0$ and the alternating sine expansion
through the negative seventh-degree term give $S(z)\le\sin(\pi z)$ on
this interval. Also
\[
 S(z)\ge z\left(a_0-\frac{b_0^3}{24}
                       -\frac{b_0^7}{322560}\right)\ge\frac95z\ge0.
\]
Expanding the four powers gives the displayed rational coefficients.
\end{proof}

Define
\[
 f(u)=\begin{cases}
 u^4(3+5u^2+6u^4),&0\le u\le5/8,\\
 15/16,&5/8<u\le1.
 \end{cases}
\]
The following strict rational lower bounds include the prefactor. The
rows use the indicated polynomial representatives up to their shared
endpoints, which do not affect the integrals.
\begin{center}
\begin{tabular}{@{}llll@{}}
\toprule
Integrand and prefactor & Interval & Degree & Lower bound\\
\midrule
$\dfrac{9a_0}{48}f(u)(1-u^2/2)S(\min(u,1-u))^3$
 &$[0,1/2]$&31&$11869/10^6$\\
same &$[1/2,5/8]$&31&$30175/10^6$\\
same &$[5/8,1]$&23&$39518/10^6$\\
\bottomrule
\end{tabular}
\end{center}
For a rational polynomial $p(u)=\sum c_ju^j$ on $[r,s]$, these values
are obtained from
$\int_r^s p(u)\,du=\sum c_j(s^{j+1}-r^{j+1})/(j+1)$.
The three rows sum to more than $163/2000>2/25$; the cubic power $S^3$
is essential in each row.

\begin{proof}[Proof of Proposition~\ref{prop:upper-volume-deficit}]
For the sine-squared probability weight defining $B$, its time variance
is $t^2(1/12-1/(2\pi^2))$.  The variance identity and \eqref{eq:8} give
\[
 \mathcal E|B|^2\ge\mathcal E|A|^2
 \left(1-\frac32t^2\left(\frac1{12}-\frac1{2\pi^2}\right)\right).
\]
Indeed, parallel identification along a geodesic and Cauchy--Schwarz
show
$\mathcal E|A(s)-A(r)|^2\le(s-r)^2\mathcal E|A'|^2$.
With $u=t/\pi$ the factor in parentheses is at least $1-u^2/2$.
For $0\le u<1$, furthermore,
\[
 f(u)\le\min\{1,\beta^2\max(1,3-\beta)\},\qquad
 \beta=\frac{u^2}{1-u^2}.
\]
At $u=1$ we give the capped comparison function its continuous value $1$.
For $u\le5/8$, comparison with $\beta^2(3-\beta)$ reduces to
$6-13u^2+6u^4\ge0$. On this interval $f$ is increasing and
$f(5/8)=7511875/8388608<1$, which also verifies the cap by $1$.  The comparison function $\min\{1,\beta^2\max(1,3-\beta)\}$, as a function of $u$, is nondecreasing and remains $1$ after it first reaches $1$; at $u=5/8$ its value is
$57500/59319>15/16$, which handles the rest of the interval.
Lemmas~\ref{lem:upper-polar} and \ref{lem:upper-radial} therefore show
that the coefficient of $s_0$ saved from the round volume is at least
\[
 \eta_+=\frac{9\pi}{48}\int_0^1
 f(u)(1-u^2/2)\sin^3(\pi u)\,du.
\]
The first three rational bounds above give $\eta_+>2/25>3/40$, proving
\eqref{eq:7}.

\end{proof}

\subsection{A moment obstruction}

The upper volume deficit is proved. Together with the characteristic
identities and the two norm gaps, it leaves a single
possible unequal-moment integer case when both blocks are nonzero, up to
exchanging the blocks. The next scalar lemma
will exclude that case. Its proof reduces a constrained probability
measure to two support points; all numerical signs are established by
explicit rational coefficients.

For $0\le z\le2$, define
\[
 H(z)=\sqrt z\left(1+\frac{z^2}{40}-\frac{z^4}{1152}\right),\qquad
 \Phi(z)=16H(z)^2+9z(1-z).
\]

\begin{lemma}[Two constrained moments]\label{lem:moment-obstruction}
Every probability measure $\mu$ on $[0,2]$ satisfying
\[
 \int z^2\,d\mu\ge1,\qquad
 \int(28z+6z^2)\,d\mu\le25
\]
satisfies
\[
 16\left(\int H\,d\mu\right)^2-\int\Phi\,d\mu<0.
\]
\end{lemma}
\begin{proof}
The feasible measures form a nonempty weakly compact convex set.
Suppose the continuous functional
$J(\mu)=16(\int H\,d\mu)^2-\int\Phi\,d\mu$ has a nonnegative maximum,
attained at $\mu$, and set $d=\int H\,d\mu$.  First $d\ge11/20$.
Indeed,
\[
 1+z^2/40-z^4/1152\ge1+(31/1440)z^2,
\]
and
\[
 25+\frac{31}{45}z^2-\frac{346}{25}z
 \ge\frac{17}{225}>0\qquad(0\le z\le2).
\]
The last quadratic is decreasing on this interval and has the stated
value at $2$.  It follows that
$\Phi(z)\ge16(11/20)^2z^2$.  The moment constraint and $J(\mu)\ge0$
now give the claim about $d$.

For another feasible measure $\nu$, direct expansion with
$\ell=32dH-\Phi$ gives
\[
 J((1-t)\mu+t\nu)-J(\mu)
 =t\left(\int\ell\,d\nu-\int\ell\,d\mu\right)
 +16t^2\left(\int H\,d\nu-d\right)^2.
\]
For $t\downarrow0$, maximality shows that $\mu$ maximizes the linear
functional $\int\ell$.  This conclusion does not require concavity of
$J$.

Put $g_1(z)=28z+6z^2-25$ and $g_2(z)=1-z^2$.
The measure with mass $9/32$ at $2$ and the rest at $0$ has respective
constraint expectations $-5/2$ and $-1/8$, so both constraints are
strictly feasible.  The compact convex set of triples
\[
 \left\{\left(\int g_1\,d\nu,\int g_2\,d\nu,\int\ell\,d\nu\right):
                 \nu\text{ a probability measure on }[0,2]\right\}
\]
is disjoint from the open set of triples with first two coordinates
negative and third coordinate greater than $\int\ell\,d\mu$.
Separation gives coefficients $-\alpha,-\gamma,\beta$, with
$\alpha,\gamma,\beta\ge0$.  The last coefficient cannot be zero:
otherwise the strictly feasible triple would violate separation.
Normalize $\beta=1$.  Evaluating on $\mu$ and on point masses shows
that $\mu$ is supported on the global maxima of
\[
 F_d(z)=32dH(z)-\Phi(z)
          -\alpha(28z+6z^2)-\gamma(1-z^2),
 \qquad \alpha,\gamma\ge0.
\]
For completeness, separation bounds
$\int(\ell-\alpha g_1-\gamma g_2)\,d\nu$ by
$\int\ell\,d\mu$ for every $\nu$; for $\mu$ the reverse inequality
follows from its nonpositive constraints.  Equality and then point
masses give precisely the asserted support condition.

We next verify that $F_d'''(z)>0$ on $(0,2]$.  Set
\[
 L(z)=\frac38+\frac{3z^2}{64}-\frac{35z^4}{1024},\qquad
 V(z)=\frac{24}{5}-\frac{16z^2}{15}-\frac{7z^4}{48}
                         +\frac{7z^6}{1152}.
\]
Direct differentiation gives $H'''(z)=L(z)z^{-5/2}$ and
$\Phi'''(z)=V(z)$.  For $x=\sqrt{z/2}$, the polynomial inequalities
\[
 L(2x^2)>0,\qquad
 D_*(x):=\frac{88}{5}L(2x^2)-\frac{17}{3}x^5V(2x^2)>0
 \qquad(0\le x\le1)
\]
are certified below.  Since $32d\ge88/5$ and $4\sqrt2<17/3$, they imply
$32dL(z)-z^{5/2}V(z)>0$ when $V(z)\ge0$; when $V(z)<0$ it is immediate.
The constraint polynomials have zero third derivative, proving the
claim.

Also $F_d'(0+)=+\infty$.  Since $F_d''$ is strictly increasing, $F_d'$
can cross from positive to negative at most once.  There is at most
one interior local maximum; an interval of maxima is impossible.
The endpoint $0$ is not a maximum, so the support contains at most
one interior point and the endpoint $2$.  A point mass is infeasible:
its first moment condition would give $z\ge1$, and then
$28z+6z^2\ge34$.  Thus the support is $\{w,2\}$, with weight $t$ at $2$.
Necessarily $0<w<1$, and the moment constraints become
\[
 t_-(w):=\frac{1-w^2}{4-w^2}\le t\le
 t_+(w):=\frac{25-28w-6w^2}{80-28w-6w^2}.
\]
The denominators are positive.  Nonemptiness reduces to
$(2-w)(10-37w)\ge0$, hence $0<w\le10/37$.

For fixed $w$, $J$ is a convex quadratic in $t$, so it is bounded above
by the larger of its two endpoint values.  Those values are negative:
\[
 16t(1-t)\bigl(H(2)-H(w)\bigr)^2
 >9\bigl((1-t)(w^2-w)+2t\bigr)
 \qquad(t=t_-(w),\ t_+(w)).
\]
To verify these inequalities exactly, put
\[
 x=\sqrt{w/2}\le\frac{37}{100},\qquad
 R_*(x)=1+\frac{x^4}{10}-\frac{x^8}{72}.
\]
For each of
\[
 (A_-,B_-)=(1-w^2,4-w^2),\qquad
 (A_+,B_+)=(25-28w-6w^2,80-28w-6w^2),\quad w=2x^2,
\]
the required inequality, after multiplication by its positive
$B_\pm^2$, is $P_\pm(x)>0$, where
\[
 P_\pm(x)=32A_\pm(B_\pm-A_\pm)
                \bigl(R_*(1)-xR_*(x)\bigr)^2
 -9\bigl((B_\pm-A_\pm)(w^2-w)+2A_\pm\bigr)B_\pm.
\]
This uses $H(2)-H(w)=\sqrt2(R_*(1)-xR_*(x))$.
The endpoint $w=10/37$, at which the weight interval can collapse, is
included in the following strict certificates.

Here are all the coefficient bounds used in this proof.  If
$p(r+(s-r)y)=\sum_{j=0}^n a_jy^j$, define
\[
 h_i=\sum_{j=0}^i a_j\frac{\binom{i}{j}}{\binom nj}
 \quad(0\le i\le n).
\]
Binomial expansion gives
$p(r+(s-r)y)=\sum_i h_i\binom ni y^i(1-y)^{n-i}$; thus positive $h_i$
give positivity on $[r,s]$.  If the original coefficients are
$p(x)=\sum p_kx^k$, they are obtained explicitly from
\[
 a_j=(s-r)^j\sum_{k=j}^n p_k\binom{k}{j}r^{k-j}.
\]
The degrees and strict lower bounds for the resulting coefficients are
\begin{center}
\begin{tabular}{@{}llll@{}}
\toprule
Polynomial & Degree & Interval & Lower bound\\
\midrule
$L(2x^2)$ &8&$[0,1]$&$1/100$\\
$D_*(x)$ &17&$[j/4,(j+1)/4]$, $j=0,1,2,3$&$2/5$\\
$P_-(x)$ &22&$[0,37/100]$&$1$\\
$P_+(x)$ &22&$[0,37/100]$&$1$\\
\bottomrule
\end{tabular}
\end{center}
For direct reproduction, the first two polynomial inputs are
\[
 L(2x^2)=\frac38+\frac3{16}x^4-\frac{35}{64}x^8,
 \qquad
 V(2x^2)=\frac{24}{5}-\frac{64}{15}x^4-\frac73x^8+\frac7{18}x^{12}.
\]
For example the least coefficient of $L(2x^2)$ is $1/64$; the least
coefficient of $D_*$ over its four intervals is
$2459668689/5793382400>2/5$.  For $P_-$ and $P_+$ the respective least
coefficients are
\[
 \begin{aligned}
 \frac{58454914732902747377800095783827783501884053}
      {50000000000000000000000000000000000000000000}&>1,\\
 \frac{132428032291164570287168821949356607977047}
      {770000000000000000000000000000000000000}&>1.
 \end{aligned}
\]
All inputs and endpoints in these coefficient calculations are rational.
This proves the derivative and endpoint signs.  The strict endpoint
inequalities say exactly that $J(t_-(w)),J(t_+(w))<0$, contradicting its
assumed nonnegative maximum.
\end{proof}

\begin{proposition}[Balanced Weyl moments]\label{prop:balanced-moments}
If neither Weyl block is identically zero, then
$\tau=0$ and $\mathbf E v^2=\mathbf E b^2$.
\end{proposition}
\begin{proof}
By \eqref{eq:5} and $v+b<2$,
$2\le s_0<4$.  Equations \eqref{eq:6} and \eqref{eq:7} give
\[
 \chi\le2(1-3s_0/40)(1+s_0)<7,
 \qquad
 \frac{|\tau|}{\chi}
 =\frac23\frac{|\mathbf E v^2-\mathbf E b^2|}{1+s_0}
 \le\frac23\frac{s_0-2}{1+s_0}<\frac4{15}.
\]
Here the first polynomial is increasing on $[2,4]$ and equals $7$ at
$4$.  The integers $\chi=2+n_++n_-$ and $\tau=n_+-n_-$ have the same
parity.  If $\tau\ne0$, the only possibility is $\chi=5$, $|\tau|=1$.
Exchange the blocks so that $P=\mathbf E b^2$ is the smaller moment.
The ratio in \eqref{eq:6} then gives
\[
 \mathbf E v^2=\frac{13P+3}{7}.
\]
Combining this with $v^2<(2-b)^2$ yields
\[
 P\ge1,\qquad28\mathbf E b+6P\le25.
\]

Integrating \eqref{eq:3} for $b$, using $\bar\sigma\le1$ and
$\bar e\ge5|\nabla b|^2/3$, gives
\[
 P-\mathbf E b\ge\frac49\mathbf E|\nabla\sqrt b|^2.
\]
The integrability of $\bar e/b$ proved with \eqref{eq:3} implies
$c=\sqrt b\in H^1$ by regularization. Put
\[
 G(c)=c+\frac{c^5}{40}-\frac{c^9}{1152},\qquad H(b)=G(c).
\]
The range of $c$ is bounded, and $G$ has bounded derivative there, so the
Sobolev chain rule applies, including at $c=0$. On $0\le b\le2$,
\[
 G'(c)=1+\frac{b^2}{8}-\frac{b^4}{128}
 \in\left[0,\frac{\sqrt{4+b^2}}2\right]
 \subseteq\left[0,\frac{\sqrt K}2\right].
\]
For the first upper bound, put $z=b^2/4\in[0,1]$ and square the positive
quantity $1+z/2-z^2/8$; its square is
$1+z-z^3/8+z^4/64\le1+z$.
The variance inequality \eqref{eq:4}, extended by Sobolev approximation,
therefore gives
\[
 \operatorname{Var}(H(b))
 \le\frac14\mathbf E|\nabla\sqrt b|^2
 \le\frac9{16}(P-\mathbf E b).
\]
Equivalently,
\[
 J:=16(\mathbf E H(b))^2-\mathbf E\Phi(b)\ge0,
 \qquad \Phi(b)=16H(b)^2+9b(1-b).
 \tag{9}\label{eq:9}
\]
The distribution of $b$ is a probability measure on $[0,2]$ satisfying
the two moment constraints.  Lemma~\ref{lem:moment-obstruction}
contradicts \eqref{eq:9}.  Thus $\tau=0$, and \eqref{eq:6} gives equality
of the two second moments.
\end{proof}
